\begin{parts}

\part[\textbf{الف)}]
برای به‌دست آوردن مدل گذار، تعداد رفتن‌ها از هر \((s,a)\) به حالت بعدی را حساب می‌کنیم و بعد تقسیم بر کل دفعات همان \((s,a)\) می‌کنیم:

\[
\hat{T}(s,a,s')=\frac{N(s,a,s')}{N(s,a)}
\]

برای حالت \(X\)، کنش \(a\) سه بار آمده است؛ دو بار به \(Y\) و یک بار به \(Z\):

\[
\hat{T}(X,a,Y)=\frac{2}{3},\qquad
\hat{T}(X,a,Z)=\frac{1}{3}
\]

برای کنش \(b\) در حالت \(X\)، هر دو بار به خود \(X\) رفته‌ایم:

\[
\hat{T}(X,b,X)=1
\]

برای حالت \(Y\)، کنش \(a\) سه بار آمده است؛ دو بار به \(Z\) و یک بار به \(Y\):

\[
\hat{T}(Y,a,Z)=\frac{2}{3},\qquad
\hat{T}(Y,a,Y)=\frac{1}{3}
\]

برای کنش \(b\) در حالت \(Y\)، فقط یک بار آمده و به \(X\) رفته است:

\[
\hat{T}(Y,b,X)=1
\]

برای حالت \(Z\)، کنش \(a\) دو بار آمده و هر دو بار در خود \(Z\) مانده است:

\[
\hat{T}(Z,a,Z)=1
\]

برای کنش \(b\) در حالت \(Z\)، فقط یک بار آمده و به \(X\) رفته است:

\[
\hat{T}(Z,b,X)=1
\]

بقیه‌ی حالت‌هایی که در داده نیامده‌اند، احتمالشان صفر در نظر گرفته می‌شود.

\part[\textbf{ب)}]
برای پاداش، برای هر حالت و کنش، میانگین پاداش‌های همان قسمت را می‌گیریم:

\[
\hat{R}(s,a,s')=\text{average reward for }(s,a,s')
\]

اینجا پاداش فقط به \(s\) و \(a\) بستگی دارد و حالت بعدی مقدارش را عوض نکرده است. پس داریم:

\[
\hat{R}(X,a,Y)=3,\qquad
\hat{R}(X,a,Z)=3
\]

\[
\hat{R}(X,b,X)=1
\]

\[
\hat{R}(Y,a,Z)=4,\qquad
\hat{R}(Y,a,Y)=4
\]

\[
\hat{R}(Y,b,X)=2
\]

\[
\hat{R}(Z,a,Z)=0
\]

\[
\hat{R}(Z,b,X)=5
\]

برای بقیه‌ی گذارها هم چون احتمالشان صفر است، لازم نیست پاداشی حساب کنیم.

\part[\textbf{ج)}]
فرمول یک مرحله از \lr{Value Iteration} این است:

\[
V_1(s)=\max_a \sum_{s'} \hat{T}(s,a,s')
\left[
\hat{R}(s,a,s')+\gamma V_0(s')
\right]
\]

چون در صورت سؤال گفته شده است که \(V_0(s)=0\)، قسمت \(\gamma V_0(s')\) صفر می‌شود. پس فقط میانگین پاداش هر کنش را حساب می‌کنیم:

\[
V_1(s)=\max_a \sum_{s'} \hat{T}(s,a,s')\hat{R}(s,a,s')
\]

برای حالت \(X\):

\[
Q_1(X,a)
=
\frac{2}{3}(3)+\frac{1}{3}(3)
=
3
\]

\[
Q_1(X,b)=1(1)=1
\]

پس:

\[
V_1(X)=\max\{3,1\}=3
\]

\[
\pi(X)=a
\]

برای حالت \(Y\):

\[
Q_1(Y,a)
=
\frac{2}{3}(4)+\frac{1}{3}(4)
=
4
\]

\[
Q_1(Y,b)=1(2)=2
\]

پس:

\[
V_1(Y)=\max\{4,2\}=4
\]

\[
\pi(Y)=a
\]

برای حالت \(Z\):

\[
Q_1(Z,a)=1(0)=0
\]

\[
Q_1(Z,b)=1(5)=5
\]

پس:

\[
V_1(Z)=\max\{0,5\}=5
\]

\[
\pi(Z)=b
\]

در نتیجه داریم:

\[
V_1(X)=3,\qquad \pi(X)=a
\]

\[
V_1(Y)=4,\qquad \pi(Y)=a
\]

\[
V_1(Z)=5,\qquad \pi(Z)=b
\]

\part[\textbf{د)}]
بله، \lr{Value Iteration} در این MDP همگرا می‌شود. دلیلش این است که تعداد حالت‌ها و کنش‌ها محدود است و مقدار ضریب تخفیف هم کمتر از ۱ است:

\[
\gamma=0.9<1
\]

یعنی اثر پاداش‌های خیلی دورتر، کم‌کم کمتر می‌شود و مقدارها به یک جواب ثابت می‌رسند. این‌که حالت \(Z\) با کنش \(a\) در خودش می‌ماند، مشکلی ایجاد نمی‌کند؛ چون باز هم ضریب تخفیف کمتر از ۱ است.

\end{parts}